% Part: incompleteness
% Chapter: incompleteness-provability
% Section: first-incompleteness-thm

\documentclass[../../../include/open-logic-section]{subfiles}

\begin{document}

\olfileid{inc}{inp}{1in}

\olsection{مبرهنة عدم الاكتمال الأولى}

تهيأ لنا الآن عرض برهان غودل الأصلي لمبرهنة عدم الاكتمال الأولى.
نأخذ $\Th{T}$ نظرية مؤكسمة حسابيًا، لغتها امتداد للغة الحساب،
ونفرض أن $\Th{T}$ تشتمل على بديهيات $\Th{Q}$. ومن هذا الفرض
يلزم، على الخصوص، أن $\Th{T}$ تمثل كل دالة وعلاقة قابلتين للحساب.

قد بيّنّا أن $\Prf[\Th{T}]({\OLId{latin-ordinary}{x}},{\OLId{latin-ordinary}{y}})$ علاقة قابلة للحساب متى اعتمدنا
ترميزًا معقولًا للصيغ والبراهين بالأعداد. ومعناها أن
$\Prf[\Th{T}]({\OLId{latin-ordinary}{x}},{\OLId{latin-ordinary}{y}})$ تصدق بالضبط حين يكون ${\OLId{latin-ordinary}{x}}$ عدد غودل
لـ!!a{derivation} للـ!!{formula} ذات عدد غودل ${\OLId{latin-ordinary}{y}}$ في
$\Th{T}$. وأثبت غودل، للنظرية التي اختارها، أمرًا أدق: هذه العلاقة
عودية بدائية، مستعينًا بقائمته المؤلفة من خمس وأربعين دالة وعلاقة.
وآخرها، ${\OLId{latin-ordinary}{x}} {\OLId{latin-looped}{B}} {\OLId{latin-ordinary}{y}}$، هي $\Prf[\Th{T}]({\OLId{latin-ordinary}{x}},{\OLId{latin-ordinary}{y}})$ لا غير بالنسبة إلى
نظريته $\Th{T}$. وينبغي أن يُقرأ اصطلاحه «عودية» هنا بمعنى ما
نسميه اليوم «عودية بدائية».

ومن قابلية $\Prf[\Th{T}]({\OLId{latin-ordinary}{x}},{\OLId{latin-ordinary}{y}})$ للحساب يلزم إمكان تمثيلها في
$\Th{T}$؛ ولتكن $\OPrf[\Th{T}]({\OLId{latin-ordinary}{x}},{\OLId{latin-ordinary}{y}})$ صيغتها الممثلة.
ونعرّف $\OProv[\Th{T}]({\OLId{latin-ordinary}{y}})$ بأنها
$\lexists[{\OLId{latin-ordinary}{x}}][\OPrf[\Th{T}]({\OLId{latin-ordinary}{x}},{\OLId{latin-ordinary}{y}})]$. فهذه هي العلاقة السادسة
والأربعون في قائمة غودل، $\fn{Bew}({\OLId{latin-ordinary}{y}})$. وقد نبه إلى أنها العلاقة
الوحيدة التي «لا يمكن الجزم بأنها عودية». ولعل مراده أن قابلية
حسابها لا تظهر من تعريفها؛ ثم دلت النتائج اللاحقة على أنها
غير قابلة للحساب بالفعل.

لتكن $\Th{T}$ نظرية !!{axiomatizable} تحتوي $\Th{Q}$.
فتكون $\Prf[\Th{T}]({\OLId{latin-ordinary}{x}},{\OLId{latin-ordinary}{y}})$ قابلة للقرار، وتمثلها في $\Th{Q}$
!!a{formula} هي $\OPrf[\Th{T}]({\OLId{latin-ordinary}{x}},{\OLId{latin-ordinary}{y}})$. ونأخذ
$\OProv[\Th{T}]({\OLId{latin-ordinary}{y}})$ كما عرّفناها. وتعطي لمّة النقطة الثابتة صيغة
$!G_\Th{T}$ بحيث !!{derive}s النظام $\Th{Q}$، ومن ثم $\Th{T}$،
التكافؤ
\begin{equation}
\ollabel{eqn:qpf}
!G_\Th{T} \liff \lnot \OProv[\Th{T}](\gn{!G_\Th{T}}).
\end{equation}
فمؤدى $!G_\Th{T}$ أن «$!G_\Th{T}$ ليست !!{derivable} في
$\Th{T}$».

\begin{lem}\ollabel{lem:cons-G-unprov}
متى كانت $\Th{T}$ متسقة و!!{axiomatizable} وتمتد $\Th{Q}$،
لزم $\Th{T}\Proves/ !G_\Th{T}$.
\end{lem}

\begin{proof}
لنفرض أن النظام $\Th{T}$ !!{derive}s $!G_\Th{T}$.
فيوجد \emph{حقًا} !!a{derivation}، وله عدد غودل ${\OLId{latin-ordinary}{m}}$ يحقق
$\Prf[\Th{T}]({\OLId{latin-ordinary}{m}},\Gn{!G_\Th{T}})$. وحينئذ النظام $\Th{Q}$
!!{derive}s الجملة $\OPrf[\Th{T}](\num {\OLId{latin-ordinary}{m}},\gn{!G_\Th{T}})$،
ومنها !!{derive}s النظام $\Th{Q}$
$\lexists[{\OLId{latin-ordinary}{x}}][\OPrf[\Th{T}]({\OLId{latin-ordinary}{x}},\gn{!G_\Th{T}})]$،
أي $\OProv[\Th{T}](\gn{!G_\Th{T}})$ بحسب التعريف.
فبمقتضى \olref{eqn:qpf}، !!{derive}s النظام $\Th{Q}$
الصيغة $\lnot !G_\Th{T}$. ولأن $\Th{T}$ تمتد $\Th{Q}$،
فإن $\Th{T}$ !!{derive}s تلك الصيغة أيضًا.
فإذا كان $\Th{T}$ !!{derive}s $!G_\Th{T}$، لزم أنه !!{derive}s
$\lnot !G_\Th{T}$ كذلك؛ وهذا ينافي الاتساق.
\end{proof}

\begin{defn}
\ollabel{thm:oconsis-q}
نصف النظرية $\Th{T}$ بأنها \emph{متسقة-$\omega$} إذا تحقق الشرط
التالي لكل جملة $\lexists[{\OLId{latin-ordinary}{x}}][!A({\OLId{latin-ordinary}{x}})]$: إن كان النظام $\Th{T}$
!!{derive}s $\lnot !A(\num {\OLMathNumeral{0}})$ و$\lnot !A(\num {\OLMathNumeral{1}})$ و$\lnot !A(\num {\OLMathNumeral{2}})$،
و\dots، امتنع أن تثبت $\Th{T}$ الجملة $\lexists[{\OLId{latin-ordinary}{x}}][!A({\OLId{latin-ordinary}{x}})]$.
\end{defn}

الاتساق-$\omega$ يستلزم الاتساق المعتاد. ولبيانه نفرض العكس:
إذا كانت $\Th{T}$ غير متسقة كان $\Th{T}\Proves !A$ لأي $!A$.
وعليه، متى كانت $\Th{T}$ غير متسقة، فإنها !!{derive}s جميع
$\lnot !A(\num {\OLId{latin-ordinary}{n}})$، لكل ${\OLId{latin-ordinary}{n}}$، وتثبت معها
$\lexists[{\OLId{latin-ordinary}{x}}][!A({\OLId{latin-ordinary}{x}})]$. فتكون $\Th{T}$ غير متسقة-$\omega$.
وبعكس النقيض، لا بد للنظرية $\Th{T}$ المتسقة-$\omega$ أن تكون
متسقة.

\begin{lem}\ollabel{lem:omega-cons-G-unref}
متى كانت $\Th{T}$ متسقة-$\omega$ و!!{axiomatizable} وتمتد
$\Th{Q}$، لزم $\Th{T}\Proves/\lnot !G_\Th{T}$.
\end{lem}

\begin{proof}
نثبت أن النظام $\Th{T}$ إذا كان !!{derive}s $\lnot !G_\Th{T}$
لم يكن متسقًا-$\omega$. فنفرض أن $\Th{T}$ !!{derive}s
$\lnot !G_\Th{T}$. إن كانت $\Th{T}$ غير متسقة، فهي غير
متسقة-$\omega$ وقد تم المطلوب. وإن كانت $\Th{T}$ متسقة، امتنع
عليها !!{derive} الصيغة $!G_\Th{T}$ بمقتضى
\olref{lem:cons-G-unprov}. وبانتفاء !!{derivation} لـ$!G_\Th{T}$
في $\Th{T}$، يكون النظام $\Th{Q}$ !!{derive}s
\[
\lnot \OPrf[\Th{T}](\num {\OLMathNumeral{0}}, \gn{!G_\Th{T}}), \lnot \OPrf[\Th{T}](\num {\OLMathNumeral{1}},
\gn{!G_\Th{T}}), \lnot \OPrf[\Th{T}](\num {\OLMathNumeral{2}}, \gn{!G_\Th{T}}), \dots
\]
وتثبتها $\Th{T}$ كذلك. ولكن $\lnot !G_\Th{T}$ مكافئة،
بحسب \olref{eqn:qpf}، لـ
$\lexists[{\OLId{latin-ordinary}{x}}][\OPrf[\Th{T}]({\OLId{latin-ordinary}{x}},\gn{!G_\Th{T}})]$.
فيجتمع في $\Th{T}$ ما يمنع اتساقها-$\omega$.
\end{proof}

\begin{prob}
  قد ثبت أن المتسقة-$\omega$ متسقة. أثبت أن العكس غير لازم،
  ببناء نظرية متسقة غير متسقة-$\omega$؛ وتحديدًا، بيّن أن
  $\Th{Q}\cup\{\lnot !G_\Th{Q}\}$ متسقة لكنها غير متسقة-$\omega$.
\end{prob}

\begin{thm}
\ollabel{thm:first-incompleteness} إذا كانت $\Th{T}$ نظرية
متسقة-$\omega$ و!!{axiomatizable} تمتد $\Th{Q}$،
فإن $\Th{T}$ غير تامة.
\end{thm}

\begin{proof}
  اتساق $\Th{T}$-$\omega$ يقتضي اتساقها، فيلزم
  $\Th{T}\Proves/ !G_\Th{T}$ من \olref{lem:cons-G-unprov}.
  وتقتضي \olref{lem:omega-cons-G-unref} أيضًا
  $\Th{T}\Proves/\lnot !G_\Th{T}$. فلا تكون $\Th{T}$ تامة:
  لا !!{derive}s $!G_\Th{T}$ ولا نفيها $\lnot !G_\Th{T}$.
\end{proof}

\end{document}
